\documentclass[12pt]{article}
\usepackage[T1]{fontenc}
\usepackage[utf8]{inputenc}
\usepackage{lmodern}
\usepackage{textcomp}
\usepackage{enumitem}
\usepackage{geometry}
\geometry{margin=1in}
\begin{document}
\begin{center}
Answer Key - Philosophy - Undergraduate Level Assessment 5 \
\small (Answers are ordered exactly as in the question paper.)
\end{center}

\section*{Multiple Choice}
\begin{enumerate}[label=\arabic*.]
\item \textbf{Answer:} C
\\ \textit{Options:} A) red herring; B) false analogy; C) appeal to force; D) appeal to tradition; E) circular argument
\\ \textit{Note:} The correct answer is "appeal to force" because this fallacy occurs when an arguer uses threats or intimidation to persuade someone to accept a conclusion rather than providing valid reasoning or evidence.
\item \textbf{Answer:} C
\\ \textit{Options:} A) has no value.; B) has value in part because it is useful to society.; C) has value solely because it is useful to society.; D) has value solely because it is useful to the agent.
\\ \textit{Note:} Hume argues that justice is a social construct that derives its value from its utility in promoting societal order and cooperation, rather than from any inherent moral property.
\item \textbf{Answer:} C
\\ \textit{Options:} A) an agent's inability to acknowledge his or her contribution to climate change; B) the fragmentation of the effects of greenhouse gas emissions across generations; C) the vast number of individuals and agents who contribute to climate change; D) the idea that it's not rational for each individual to restrict his or her own pollution
\\ \textit{Note:} The correct answer highlights that "the fragmentation of agency" signifies the diverse array of individuals and entities whose collective actions and decisions contribute to the complexity and challenge of addressing climate change, reflecting a key philosophical concern regarding moral responsibility and accountability in environmental ethics.
\item \textbf{Answer:} D
\\ \textit{Options:} A) justice; B) fairness; C) respect for other persons; D) all of the above
\\ \textit{Note:} Arthur believes that values such as justice, equality, and community are more fundamental than individual rights and desert, as they provide the foundational principles that guide and shape ethical considerations and social interactions.
\item \textbf{Answer:} B
\\ \textit{Options:} A) Epistemology; B) Morality; C) Religion; D) Aesthetics
\\ \textit{Note:} In his final work, Laws, Plato emphasizes the importance of moral principles and ethical governance, indicating a significant shift from his earlier focus on cosmological theories to the practical implications of morality in society.
\end{enumerate}

\section*{True / False}
\begin{enumerate}[label=\arabic*.]
\setcounter{enumi}{5}
\item \textbf{Answer:} False
\\ \textit{Options:} A) True; B) False
\\ \textit{Note:} The Oral Torah consists of teachings and interpretations passed down orally and is not written in the Tanakh, which refers specifically to the written texts of the Hebrew Bible.
\item \textbf{Answer:} False
\\ \textit{Options:} A) True; B) False
\\ \textit{Note:} Stevenson's primary aim in the paper is to explore the nature of moral judgments and their implications rather than to provide a definitive account of what makes right actions right.
\item \textbf{Answer:} True
\\ \textit{Options:} A) True; B) False
\\ \textit{Note:} The Digambara sect of Jainism allows women to lead a life of semi-renunciation because they believe that, while ultimate liberation is not attainable for women in their present form, they can still engage in spiritual practices and lead a lifestyle that approximates renunciation, albeit with certain restrictions.
\item \textbf{Answer:} True
\\ \textit{Options:} A) True; B) False
\\ \textit{Note:} In a conditional syllogism, for the argument to be valid, the minor premise must affirm the antecedent or deny the consequent, ensuring that the logical structure aligns with the truth conditions of the conditional statement.
\item \textbf{Answer:} False
\\ \textit{Options:} A) True; B) False
\\ \textit{Note:} Anscombe contends that the concept of moral obligations is problematic and lacks coherence without a framework of divine command or a clear understanding of human purpose, leading her to reject the idea that moral obligations can be entirely coherent and well-defined.
\end{enumerate}

\section*{Long Form Response}
\begin{enumerate}[label=\arabic*.]
\setcounter{enumi}{10}
\item \textbf{Answer:} To evaluate the argument consisting of (A ∨ B) ⊃ C, C ⊃ \textasciitilde{}D, and D ⊃ A using indirect truth tables, we find that the argument is invalid. A counterexample occurs when A, B, and C are all false while D is true. In this scenario, (A ∨ B) is false, making (A ∨ B) ⊃ C true by vacuous truth; however, C being false leads to C ⊃ \textasciitilde{}D being false, as D is true. This demonstrates that not all premises can simultaneously hold while maintaining the conclusion, revealing the argument's inconsistency and underscoring the potential for flawed reasoning in philosophical discourse.
\\ \textit{Note:} correct because it demonstrates the invalidity of the argument by providing a specific counterexample that reveals a scenario where the premises are true while the conclusion fails, thus highlighting the logical inconsistency of the argument.
\item \textbf{Answer:} Moral cosmopolitanism posits that individuals have ethical obligations to assist those in need and advocate for basic human rights regardless of national boundaries. This perspective emphasizes that all human beings possess inherent dignity and worth, which creates a responsibility for individuals to transcend parochial loyalties and respond to global injustices. By promoting human rights and aiding those suffering from poverty, conflict, or oppression, individuals contribute to a more equitable world. Furthermore, moral cosmopolitanism challenges the notion that moral duties are confined to one's compatriots, urging a more inclusive approach to ethics that recognizes the interconnectedness of humanity. Ultimately, it calls for individuals to engage actively in global solidarity and humanitarian efforts, thereby fulfilling their moral obligations on a broader scale.
\\ \textit{Note:} correct because it effectively captures the essence of moral cosmopolitanism, which asserts that individuals have a moral duty to support global human rights and assist those in need, highlighting the universal dignity of all people and the importance of addressing injustices that extend beyond national borders.
\end{enumerate}

\vfill
\noindent\textit{End of Answer Key}
\end{document}